\nonstopmode{}
\documentclass[letterpaper]{book}
\usepackage[times,inconsolata,hyper]{Rd}
\usepackage{makeidx}
\makeatletter\@ifl@t@r\fmtversion{2018/04/01}{}{\usepackage[utf8]{inputenc}}\makeatother
% \usepackage{graphicx} % @USE GRAPHICX@
\makeindex{}
\begin{document}
\chapter*{}
\begin{center}
{\textbf{\huge Package `fspls2'}}
\par\bigskip{\large \today}
\end{center}
\ifthenelse{\boolean{Rd@use@hyper}}{\hypersetup{pdftitle = {fspls2: A Package to Find Minimal Transcriptional Signatures}}}{}
\ifthenelse{\boolean{Rd@use@hyper}}{\hypersetup{pdfauthor = {Lachlan Coin}}}{}
\begin{description}
\raggedright{}
\item[Type]\AsIs{Package}
\item[Title]\AsIs{A Package to Find Minimal Transcriptional Signatures}
\item[Version]\AsIs{0.1}
\item[Date]\AsIs{2026-06-25}
\item[Author]\AsIs{Lachlan Coin [aut, cre]}
\item[Imports]\AsIs{R6, jsonlite, Matrix, glmnet, tidyr, MASS, ggplot2, pROC,
binom, confintr, DBI, RSQLite, methods,tibble, ggrepel,
RColorBrewer, parallel}
\item[Maintainer]\AsIs{Lachlan Coin }\email{l.coin@imb.uq.edu.au}\AsIs{}
\item[Description]\AsIs{Identifies minimal gene signatures from multiomic data.}
\item[License]\AsIs{GPL-2}
\item[LazyLoad]\AsIs{yes}
\item[NeedsCompilation]\AsIs{no}
\item[Repository]\AsIs{CRAN}
\item[Date/Publication]\AsIs{2026-06-25 13:50:02 UTC}
\item[Config/roxygen2/version]\AsIs{8.0.0}
\item[Encoding]\AsIs{UTF-8}
\item[Suggests]\AsIs{knitr, rmarkdown}
\item[VignetteBuilder]\AsIs{knitr}
\end{description}
\Rdcontents{Contents}
\HeaderA{analysisBase}{analysisBase}{analysisBase}
%
\begin{Description}
A class that encapsulates a dataset
\end{Description}
%
\begin{Section}{Methods}
%
\begin{SubSection}{Public methods}
\begin{itemize}

\item{} \Rhref{\#method-analysisBase-initialize}{\code{analysisBase\$new()}}
\item{} \Rhref{\#method-analysisBase-name}{\code{analysisBase\$name()}}
\item{} \Rhref{\#method-analysisBase-clone}{\code{analysisBase\$clone()}}

\end{itemize}

\end{SubSection}



\hypertarget{method-analysisBase-initialize}{}
%
\begin{SubSection}{\code{analysisBase\$new()}}
Create a new instance of base class
%
\begin{SubSubSection}{Usage}

\begin{alltt}analysisBase$new(nme, dims = NULL, flags = list(), dbDir = tempdir())\end{alltt}



\end{SubSubSection}

%
\begin{SubSubSection}{Arguments}

\begin{description}

\item[\code{nme}] name
\item[\code{dims}] dims
\item[\code{flags}] a list object specifying options
\item[\code{dbDir}] location for database to be written, default to tempdir

\end{description}



\end{SubSubSection}

\end{SubSection}




\hypertarget{method-analysisBase-name}{}
%
\begin{SubSection}{\code{analysisBase\$name()}}
get dataset name;
%
\begin{SubSubSection}{Usage}

\begin{alltt}analysisBase$name()\end{alltt}



\end{SubSubSection}

%
\begin{SubSubSection}{Returns}
assigned name of dataset

\end{SubSubSection}

\end{SubSection}




\hypertarget{method-analysisBase-clone}{}
%
\begin{SubSection}{\code{analysisBase\$clone()}}
The objects of this class are cloneable with this method.
%
\begin{SubSubSection}{Usage}

\begin{alltt}analysisBase$clone(deep = FALSE)\end{alltt}



\end{SubSubSection}

%
\begin{SubSubSection}{Arguments}

\begin{description}

\item[\code{deep}] Whether to make a deep clone.

\end{description}



\end{SubSubSection}

\end{SubSection}


\end{Section}
\HeaderA{analysisEnv}{analysisEnv}{analysisEnv}
%
\begin{Description}
A class that encapsulates a FSPLS analysis of multiple datasets
\end{Description}
%
\begin{Section}{Super class}
\code{\LinkA{analysisBase}{analysisBase}} -> \code{analysisEnv}
\end{Section}
%
\begin{Section}{Methods}
%
\begin{SubSection}{Public methods}
\begin{itemize}

\item{} \Rhref{\#method-analysisEnv-initialize}{\code{analysisEnv\$new()}}
\item{} \Rhref{\#method-analysisEnv-getTodo}{\code{analysisEnv\$getTodo()}}
\item{} \Rhref{\#method-analysisEnv-clear_db}{\code{analysisEnv\$clear\_db()}}
\item{} \Rhref{\#method-analysisEnv-savePvalsAndNextVars}{\code{analysisEnv\$savePvalsAndNextVars()}}
\item{} \Rhref{\#method-analysisEnv-select}{\code{analysisEnv\$select()}}
\item{} \Rhref{\#method-analysisEnv-clone}{\code{analysisEnv\$clone()}}

\end{itemize}

\end{SubSection}




\hypertarget{method-analysisEnv-initialize}{}
%
\begin{SubSection}{\code{analysisEnv\$new()}}
Create a new instance
%
\begin{SubSubSection}{Usage}

\begin{alltt}analysisEnv$new(flags = list(), dbDir = tempdir())\end{alltt}



\end{SubSubSection}

%
\begin{SubSubSection}{Arguments}

\begin{description}

\item[\code{flags}] a list object specifying options
\item[\code{dbDir}] location for database to be written. If NULL then no DB 
Get todo

\end{description}



\end{SubSubSection}

\end{SubSection}




\hypertarget{method-analysisEnv-getTodo}{}
%
\begin{SubSection}{\code{analysisEnv\$getTodo()}}
%
\begin{SubSubSection}{Usage}

\begin{alltt}analysisEnv$getTodo(flags, phens)\end{alltt}



\end{SubSubSection}

%
\begin{SubSubSection}{Arguments}

\begin{description}

\item[\code{flags}] list of options
\item[\code{phens}] phenotypes

\end{description}



\end{SubSubSection}

%
\begin{SubSubSection}{Returns}
object outlining what is left to do for next iteration, 
Clear the databases.  Only necessary if using a database and re-running analyses

\end{SubSubSection}

\end{SubSection}




\hypertarget{method-analysisEnv-clear_db}{}
%
\begin{SubSection}{\code{analysisEnv\$clear\_db()}}
%
\begin{SubSubSection}{Usage}

\begin{alltt}analysisEnv$clear_db(drop = F, exclude = "vars", datasH = NULL)\end{alltt}



\end{SubSubSection}

%
\begin{SubSubSection}{Arguments}

\begin{description}

\item[\code{drop}] completely drop tables and start from scratch
\item[\code{exclude}] tables to exclue from clearing
\item[\code{datasH}] list of dataH objects to clear
Get the next vars in iteration.  Internal function

\end{description}



\end{SubSubSection}

\end{SubSection}




\hypertarget{method-analysisEnv-savePvalsAndNextVars}{}
%
\begin{SubSection}{\code{analysisEnv\$savePvalsAndNextVars()}}
%
\begin{SubSubSection}{Usage}

\begin{alltt}analysisEnv$savePvalsAndNextVars(
  flags,
  phens,
  vars_l_todo,
  comb2,
  data_nme,
  k1
)\end{alltt}



\end{SubSubSection}

%
\begin{SubSubSection}{Arguments}

\begin{description}

\item[\code{flags}] list of options
\item[\code{phens}] phenotypes
\item[\code{vars\_l\_todo}] vars\_l\_todo
\item[\code{comb2}] results
\item[\code{data\_nme}] name of dataset
\item[\code{k1}] which repetition

\end{description}



\end{SubSubSection}

%
\begin{SubSubSection}{Returns}
object outlining what is left to do
main function for variable selection

\end{SubSubSection}

\end{SubSection}




\hypertarget{method-analysisEnv-select}{}
%
\begin{SubSection}{\code{analysisEnv\$select()}}
%
\begin{SubSubSection}{Usage}

\begin{alltt}analysisEnv$select(
  datasH,
  flags,
  transform_y,
  phens = datasH[[1]]$pheno()$all,
  useDB = F
)\end{alltt}



\end{SubSubSection}

%
\begin{SubSubSection}{Arguments}

\begin{description}

\item[\code{datasH}] a list of dataH objects
\item[\code{flags}] list of options
\item[\code{transform\_y}] transformation object
\item[\code{phens}] list of phenotypes
\item[\code{useDB}] boolean to indicate if results should be saved to database

\end{description}



\end{SubSubSection}

\end{SubSection}




\hypertarget{method-analysisEnv-clone}{}
%
\begin{SubSection}{\code{analysisEnv\$clone()}}
The objects of this class are cloneable with this method.
%
\begin{SubSubSection}{Usage}

\begin{alltt}analysisEnv$clone(deep = FALSE)\end{alltt}



\end{SubSubSection}

%
\begin{SubSubSection}{Arguments}

\begin{description}

\item[\code{deep}] Whether to make a deep clone.

\end{description}



\end{SubSubSection}

\end{SubSection}


\end{Section}
\HeaderA{check\_flags}{Plot the results from evaluateAll}{check.Rul.flags}
%
\begin{Description}
Plot the results from evaluateAll
\end{Description}
%
\begin{Usage}
\begin{verbatim}
check_flags(flags)
\end{verbatim}
\end{Usage}
%
\begin{Arguments}
\begin{ldescription}
\item[\code{flags}] a list of options
\end{ldescription}
\end{Arguments}
%
\begin{Value}
corrected flags
\end{Value}
\HeaderA{dataH}{dataH}{dataH}
%
\begin{Description}
A class that encapsulates a single dataset
\end{Description}
%
\begin{Section}{Super class}
\code{\LinkA{analysisBase}{analysisBase}} -> \code{dataH}
\end{Section}
%
\begin{Section}{Methods}
%
\begin{SubSection}{Public methods}
\begin{itemize}

\item{} \Rhref{\#method-dataH-initialize}{\code{dataH\$new()}}
\item{} \Rhref{\#method-dataH-multiAnglesAndPv}{\code{dataH\$multiAnglesAndPv()}}
\item{} \Rhref{\#method-dataH-clear_db}{\code{dataH\$clear\_db()}}
\item{} \Rhref{\#method-dataH-data_types}{\code{dataH\$data\_types()}}
\item{} \Rhref{\#method-dataH-split}{\code{dataH\$split()}}
\item{} \Rhref{\#method-dataH-pheno}{\code{dataH\$pheno()}}
\item{} \Rhref{\#method-dataH-dims}{\code{dataH\$dims()}}
\item{} \Rhref{\#method-dataH-nreps}{\code{dataH\$nreps()}}
\item{} \Rhref{\#method-dataH-update}{\code{dataH\$update()}}
\item{} \Rhref{\#method-dataH-select}{\code{dataH\$select()}}
\item{} \Rhref{\#method-dataH-extractPredictions}{\code{dataH\$extractPredictions()}}
\item{} \Rhref{\#method-dataH-plotData}{\code{dataH\$plotData()}}
\item{} \Rhref{\#method-dataH-getProjectedData}{\code{dataH\$getProjectedData()}}
\item{} \Rhref{\#method-dataH-getVariance}{\code{dataH\$getVariance()}}
\item{} \Rhref{\#method-dataH-makeAllModels}{\code{dataH\$makeAllModels()}}
\item{} \Rhref{\#method-dataH-evaluateAllModels}{\code{dataH\$evaluateAllModels()}}
\item{} \Rhref{\#method-dataH-clone}{\code{dataH\$clone()}}

\end{itemize}

\end{SubSection}




\hypertarget{method-dataH-initialize}{}
%
\begin{SubSection}{\code{dataH\$new()}}
Create a new instance
%
\begin{SubSubSection}{Usage}

\begin{alltt}dataH$new(
  data,
  y,
  nme,
  flags,
  transform_y = getYTransform(pow = 1, n_random = 1, perm = F),
  family = getFamily(y),
  dbDir = tempdir()
)\end{alltt}



\end{SubSubSection}

%
\begin{SubSubSection}{Arguments}

\begin{description}

\item[\code{data}] a matrix, with columns as variables and rows as samples
\item[\code{y}] phenotype matrix, with columns as outcomes and rows as samples
\item[\code{nme}] The name of the data objetct.
\item[\code{flags}] list of options, described in the vignette
\item[\code{transform\_y}] a transformation object from the function getYTransformation
\item[\code{family}] the statistical family of phenotype y, can be calculated by getFamily
\item[\code{dbDir}] dir for database to store results to speed up re-reruns.  Can be NULL, if not required

\end{description}



\end{SubSubSection}

\end{SubSection}




\hypertarget{method-dataH-multiAnglesAndPv}{}
%
\begin{SubSection}{\code{dataH\$multiAnglesAndPv()}}
Calculate the angles and pv across multiple phenotypes.  This is an internal function and should not need to be called by user
%
\begin{SubSubSection}{Usage}

\begin{alltt}dataH$multiAnglesAndPv(comb20, phens1, k1, flags, expt_id, vars_l_todo)\end{alltt}



\end{SubSubSection}

%
\begin{SubSubSection}{Arguments}

\begin{description}

\item[\code{comb20}] values from previous iteration
\item[\code{phens1}] phenotypes being used
\item[\code{k1}] k1 is the current fold
\item[\code{flags}] list of options
\item[\code{expt\_id}] and ID for the experiment (can be 0)
\item[\code{vars\_l\_todo}] the variables under consideration
clear database. This is only used if you want to clear the database.  Internal function

\end{description}



\end{SubSubSection}

\end{SubSection}




\hypertarget{method-dataH-clear_db}{}
%
\begin{SubSection}{\code{dataH\$clear\_db()}}
%
\begin{SubSubSection}{Usage}

\begin{alltt}dataH$clear_db(drop = FALSE, exclude = "vars")\end{alltt}



\end{SubSubSection}

%
\begin{SubSubSection}{Arguments}

\begin{description}

\item[\code{drop}] drop the tables instead of clear
\item[\code{exclude}] which tables not to clear
get the data types

\end{description}



\end{SubSubSection}

\end{SubSection}




\hypertarget{method-dataH-data_types}{}
%
\begin{SubSection}{\code{dataH\$data\_types()}}
%
\begin{SubSubSection}{Usage}

\begin{alltt}dataH$data_types()\end{alltt}



\end{SubSubSection}

%
\begin{SubSubSection}{Returns}
which data types are in the dataset

\end{SubSubSection}

\end{SubSection}




\hypertarget{method-dataH-split}{}
%
\begin{SubSection}{\code{dataH\$split()}}
split dataset into smaller datasets
%
\begin{SubSubSection}{Usage}

\begin{alltt}dataH$split(proportions = c(0.5, 0.5))\end{alltt}



\end{SubSubSection}

%
\begin{SubSubSection}{Arguments}

\begin{description}

\item[\code{proportions}] what proportions to split into

\end{description}



\end{SubSubSection}

%
\begin{SubSubSection}{Returns}
a list of dataH objects with data partiioned according to proportions

\end{SubSubSection}

\end{SubSection}




\hypertarget{method-dataH-pheno}{}
%
\begin{SubSection}{\code{dataH\$pheno()}}
get the phenotyeps
%
\begin{SubSubSection}{Usage}

\begin{alltt}dataH$pheno()\end{alltt}



\end{SubSubSection}

%
\begin{SubSubSection}{Returns}
phenotypes

\end{SubSubSection}

\end{SubSection}




\hypertarget{method-dataH-dims}{}
%
\begin{SubSection}{\code{dataH\$dims()}}
get the dimensions
%
\begin{SubSubSection}{Usage}

\begin{alltt}dataH$dims()\end{alltt}



\end{SubSubSection}

%
\begin{SubSubSection}{Returns}
dimensions
get the nreps

\end{SubSubSection}

\end{SubSection}




\hypertarget{method-dataH-nreps}{}
%
\begin{SubSection}{\code{dataH\$nreps()}}
%
\begin{SubSubSection}{Usage}

\begin{alltt}dataH$nreps()\end{alltt}



\end{SubSubSection}

%
\begin{SubSubSection}{Returns}
nreps

\end{SubSubSection}

\end{SubSection}




\hypertarget{method-dataH-update}{}
%
\begin{SubSection}{\code{dataH\$update()}}
update the phenotypes without remaking the entire object. You can provide many phenotypes in the initialisation stage, but only consider a subset in model fitting stage in this way.
%
\begin{SubSubSection}{Usage}

\begin{alltt}dataH$update(phens, flags, transform_y = fromJSON(flags$transform_y))\end{alltt}



\end{SubSubSection}

%
\begin{SubSubSection}{Arguments}

\begin{description}

\item[\code{phens}] phenotypes
\item[\code{flags}] list of options
\item[\code{transform\_y}] transformation object

\end{description}



\end{SubSubSection}

\end{SubSection}




\hypertarget{method-dataH-select}{}
%
\begin{SubSection}{\code{dataH\$select()}}
main function for variable selection
%
\begin{SubSubSection}{Usage}

\begin{alltt}dataH$select(analysis, phens, flags, transform_y, useDB = F)\end{alltt}



\end{SubSubSection}

%
\begin{SubSubSection}{Arguments}

\begin{description}

\item[\code{analysis}] an analysisEnv object
\item[\code{phens}] list of phenotyps
\item[\code{flags}] list of options
\item[\code{transform\_y}] transformation object
\item[\code{useDB}] boolean to indicate if results should be saved to database

\end{description}



\end{SubSubSection}

\end{SubSection}




\hypertarget{method-dataH-extractPredictions}{}
%
\begin{SubSection}{\code{dataH\$extractPredictions()}}
extract the predictions for the fitted models
%
\begin{SubSubSection}{Usage}

\begin{alltt}dataH$extractPredictions(
  all_modelsh,
  phens = all_modelsh$phens,
  flags = all_modelsh$flags,
  CV = FALSE,
  liab = T
)\end{alltt}



\end{SubSubSection}

%
\begin{SubSubSection}{Arguments}

\begin{description}

\item[\code{all\_modelsh}] fitted models from makeAllModels
\item[\code{phens}] list of phenotyps
\item[\code{flags}] list of options
\item[\code{CV}] cross validation predictions
\item[\code{liab}] return liability score, or probability (for binomial, ordinal multinomial)

\end{description}



\end{SubSubSection}

%
\begin{SubSubSection}{Returns}
a table with results

\end{SubSubSection}

\end{SubSection}




\hypertarget{method-dataH-plotData}{}
%
\begin{SubSection}{\code{dataH\$plotData()}}
plot all data for selected variables
%
\begin{SubSubSection}{Usage}

\begin{alltt}dataH$plotData(
  vars_all,
  phens = vars_all$phens,
  all_types = FALSE,
  transform_x = NULL,
  violin = FALSE,
  assoc = F
)\end{alltt}



\end{SubSubSection}

%
\begin{SubSubSection}{Arguments}

\begin{description}

\item[\code{vars\_all}] variables
\item[\code{phens}] list of phenotyps
\item[\code{all\_types}] use all types?
\item[\code{transform\_x}] which transform to use on x
\item[\code{violin}] violin plots
\item[\code{assoc}] use association

\end{description}



\end{SubSubSection}

%
\begin{SubSubSection}{Returns}
a table with results

\end{SubSubSection}

\end{SubSection}




\hypertarget{method-dataH-getProjectedData}{}
%
\begin{SubSection}{\code{dataH\$getProjectedData()}}
get data after projection
%
\begin{SubSubSection}{Usage}

\begin{alltt}dataH$getProjectedData(varnames)\end{alltt}



\end{SubSubSection}

%
\begin{SubSubSection}{Arguments}

\begin{description}

\item[\code{varnames}] varnames

\end{description}



\end{SubSubSection}

%
\begin{SubSubSection}{Returns}
projected data after projecting out varnames

\end{SubSubSection}

\end{SubSection}




\hypertarget{method-dataH-getVariance}{}
%
\begin{SubSection}{\code{dataH\$getVariance()}}
get variance of data
%
\begin{SubSubSection}{Usage}

\begin{alltt}dataH$getVariance(varnames)\end{alltt}



\end{SubSubSection}

%
\begin{SubSubSection}{Arguments}

\begin{description}

\item[\code{varnames}] varnames

\end{description}



\end{SubSubSection}

%
\begin{SubSubSection}{Returns}
variance

\end{SubSubSection}

\end{SubSection}




\hypertarget{method-dataH-makeAllModels}{}
%
\begin{SubSection}{\code{dataH\$makeAllModels()}}
fit models based on variables
%
\begin{SubSubSection}{Usage}

\begin{alltt}dataH$makeAllModels(
  vars_all,
  phens = vars_all[[1]]$phens,
  flags = vars_all[[1]]$flags,
  transform_y = fromJSON(flags$transform_y),
  useDB = F
)\end{alltt}



\end{SubSubSection}

%
\begin{SubSubSection}{Arguments}

\begin{description}

\item[\code{vars\_all}] list of variables selected by select method
\item[\code{phens}] list of phenotypes
\item[\code{flags}] flags
\item[\code{transform\_y}] an object to determine the y transformaton
\item[\code{useDB}] whether to use the DB to save the reuslt

\end{description}



\end{SubSubSection}

%
\begin{SubSubSection}{Returns}
fitted models

\end{SubSubSection}

\end{SubSection}




\hypertarget{method-dataH-evaluateAllModels}{}
%
\begin{SubSection}{\code{dataH\$evaluateAllModels()}}
evaluate the fit of models
%
\begin{SubSubSection}{Usage}

\begin{alltt}dataH$evaluateAllModels(
  all_modelsh,
  phens = all_modelsh$phens,
  flags = all_modelsh$flags,
  useDB = F
)\end{alltt}



\end{SubSubSection}

%
\begin{SubSubSection}{Arguments}

\begin{description}

\item[\code{all\_modelsh}] models fitted from makeAllModels
\item[\code{phens}] list of phenotypes
\item[\code{flags}] flags
\item[\code{useDB}] whether to save results to DB

\end{description}



\end{SubSubSection}

%
\begin{SubSubSection}{Returns}
evaluation of models

\end{SubSubSection}

\end{SubSection}




\hypertarget{method-dataH-clone}{}
%
\begin{SubSection}{\code{dataH\$clone()}}
The objects of this class are cloneable with this method.
%
\begin{SubSubSection}{Usage}

\begin{alltt}dataH$clone(deep = FALSE)\end{alltt}



\end{SubSubSection}

%
\begin{SubSubSection}{Arguments}

\begin{description}

\item[\code{deep}] Whether to make a deep clone.

\end{description}



\end{SubSubSection}

\end{SubSection}


\end{Section}
\HeaderA{getFamily}{Infer the statistical families of the phenotypes}{getFamily}
%
\begin{Description}
Infer the statistical families of the phenotypes
\end{Description}
%
\begin{Usage}
\begin{verbatim}
getFamily(y_mat, max_ordinal = getOption("max_ordinal", 10))
\end{verbatim}
\end{Usage}
%
\begin{Arguments}
\begin{ldescription}
\item[\code{y\_mat}] a matrix of phenotypes, with columns as variables and rows as samples

\item[\code{max\_ordinal}] maximum number of values before we consider an ordinal value as a continuous value
\end{ldescription}
\end{Arguments}
%
\begin{Value}
a list with the statistical families detected
\end{Value}
\HeaderA{getYTransform}{Get object for transforming the y variables}{getYTransform}
%
\begin{Description}
Get object for transforming the y variables
\end{Description}
%
\begin{Usage}
\begin{verbatim}
getYTransform(
  pows = c(1),
  offset = 0,
  n_random = 1,
  perm = F,
  norm = 1,
  CHECK = F
)
\end{verbatim}
\end{Usage}
%
\begin{Arguments}
\begin{ldescription}
\item[\code{pows}] power to raise to

\item[\code{offset}] offset to subtract

\item[\code{n\_random}] how many random variables

\item[\code{perm}] whether random is permutation

\item[\code{norm}] rescaling factor

\item[\code{CHECK}] check whether inverse function works
\end{ldescription}
\end{Arguments}
%
\begin{Value}
ggplot2 plot
\end{Value}
\HeaderA{plotEval}{Plot the results from evaluateAll}{plotEval}
%
\begin{Description}
Plot the results from evaluateAll
\end{Description}
%
\begin{Usage}
\begin{verbatim}
plotEval(
  eval3,
  shape_color = c("pheno", "subpheno"),
  shape = shape_color,
  text = "variable",
  dotsize = "nsamps",
  color = shape_color,
  linetype = shape,
  showranges = TRUE,
  txtsize = 1,
  logy = F,
  legend = F,
  sep_by = "",
  scales = "free",
  point = TRUE,
  line = TRUE,
  labelsize = 2,
  grid0 = c("cohort", "measure"),
  grid1 = "cv_full",
  title = "",
  title1 = ""
)
\end{verbatim}
\end{Usage}
%
\begin{Arguments}
\begin{ldescription}
\item[\code{eval3}] a tibble or data frame from evaluateAll method from dataH

\item[\code{shape\_color}] a list of column names to encode as both shape and color

\item[\code{shape}] defaults to shape\_color but can be specified separately

\item[\code{text}] a list of column names to include as text labels

\item[\code{dotsize}] a numerical column to encode dotsize, or a number for constant sizes

\item[\code{color}] defaults to shape\_color but can be specified separately

\item[\code{linetype}] column controlling the type of line

\item[\code{showranges}] whether to show 95\% CI as ribbon

\item[\code{txtsize}] size of text

\item[\code{logy}] whether to log y axis

\item[\code{legend}] whether to show a legend

\item[\code{sep\_by}] list of column names to separate into diff plots

\item[\code{scales}] used in faceting can be free or fixed

\item[\code{point}] whether to show dots

\item[\code{line}] whether to include line

\item[\code{labelsize}] the size of the text labels

\item[\code{grid0}] how to facet on y

\item[\code{grid1}] how to facet on x

\item[\code{title}] the title

\item[\code{title1}] columns to include in the title
\end{ldescription}
\end{Arguments}
%
\begin{Value}
ggplot2 plot
\end{Value}
\printindex{}
\end{document}
